\documentclass[a4paper, 12pt]{article}

\input{preamble.tex}

% Override Russian captions for figures, tables, contents
\renewcommand{\figurename}{Fig.}
\renewcommand{\tablename}{Table}
\renewcommand{\contentsname}{Contents}

\lstset{
  language=Python,
  basicstyle=\ttfamily\small,
  numbers=left,
  numberstyle=\tiny,
  frame=single,
  breaklines=true,
  showstringspaces=false,
  captionpos=b,
  aboveskip=3mm,
  belowskip=3mm
}

\hypersetup{
  colorlinks=true,
  linkcolor=adblue,
  filecolor=adblue,
  urlcolor=adblue,
  citecolor=adblue,
}

\begin{document}

% =============================================
% TITLE PAGE
% =============================================
\begin{titlepage}
\centering
\vspace{0.75cm}
{\Huge\bfseries Development of ADA4870 Simulator\par}
\vspace{0.25cm}
{\Large for the Output Channel of a Three-Phase High-Frequency Generator\par}
\vspace{0.25cm}
{\large (with ADA4807-2 Preamplifiers)\par}
\vspace{1.5cm}
{\large Technical Note TN-2024-01\par}
\vspace{2cm}
\vfill
\begin{flushright}
Prepared by: grandfatherpikhto\\
Date: \today
\end{flushright}
\vspace{1cm}
\end{titlepage}

\newpage
\tableofcontents
\newpage

% =============================================
% 1. INTRODUCTION
% =============================================
\section{Introduction and Objectives}

\subsection{Purpose of the Circuit}

\par This is an engineering analysis of the schematic \texttt{ada4807\_4870.asc} ---
a simulation of the output stage based on the ADA4870 operational amplifier
with ADA4807-2 current preamplifiers. The circuit is designed to amplify the
current-output signal from the AD9705 DAC and drive a low-impedance load.

Key characteristics of the ADA4870 amplifier \cite{ADA4870_ds}:
\begin{itemize}
  \item Unity-gain bandwidth: 52~MHz (large signal),
        70~MHz (small signal)
  \item Slew rate: 2500~V/µs
  \item Maximum output current: 1~A
  \item Supply voltage: from 10~V to 40~V (single or dual supply)
  \item Quiescent current: 32.5~mA (typical)
\end{itemize}

Key characteristics of the ADA4807-2 \cite{ADA4807_ds}:
\begin{itemize}
  \item Low-noise rail-to-rail operational amplifier
  \item Bandwidth 180~MHz (gain +1)
  \item Slew rate 225~V/µs
  \item Low noise: 3.1 nV/√Hz
  \item Supply range: 2.7~V … 10.5~V
\end{itemize}

\subsection{Project File Structure}

\textbf{Updated:} after refactoring, the project is structured as the package
\texttt{adasim/} with separation into \texttt{core/} (pure component calculation),
\texttt{ltspice\_io/} (parsing \texttt{.asc}/\texttt{.asy}, running LTspice),
and \texttt{report/} (plots and text report). The full structure and purpose of
each module are described in the project's \texttt{README.md} (section "Architecture"),
and are not duplicated here to avoid divergence with future changes. Key files:
\begin{itemize}
  \item \texttt{ada4807\_4870.asc} --- LTspice simulation schematic
  \item \texttt{config.yaml} --- configuration (was \texttt{config.json} in earlier versions)
  \item \texttt{out/report.txt}, \texttt{out/*.png} --- text report and plots from the run
  \item \texttt{net/*\_readable.txt} --- readable circuit report with pin names (generated directly from \texttt{.asc}, not via intermediate \texttt{.net})
\end{itemize}

\par The circuit looks as follows:

\begin{figure}[H]
\centering
\begin{circuitikz}[american, scale=0.6, transform shape]
    %%%%%%%%%%%%%%%%%%%%%%%%%%
    % Upper OpAmp (Channel A)
    %%%%%%%%%%%%%%%%%%%%%%%%%
    \draw (0,14) node[op amp] (opampA) {}
        % 1. Supply and ground
        (opampA.+) -- ++(-0.5,0) node[ground] {} 
        (opampA.up) -- ++(0,0.3) node[above] {+2.5V}
        (opampA.down) -- ++(0,-0.3) node[below] {-2.5V}
        
        % 2. Inverting input -> Point iouta -> Current source
        (opampA.-) to[short, -o] ++(-2,0) coordinate (iouta) node[left] {\(I_{outA}\)}
        (iouta) -- ++(0,-0.5) to[isource, l=$I_A$] ++(0,-1.5) node[ground]{}
        
        % 3. Back to inverting input and feedback
        (opampA.-) -- ++(0,2) coordinate (fbA) 
        to[R, l={\(R_{fa}\)}, a=249R] (fbA -| opampA.out) -- (opampA.out)
        
        % 4. Output
        (opampA.out) to[R, l={\(R_{b}\)}, a=100R, -o] ++(2.5,0) node[right] {\(OUT_A\)}; 
    
    % Model label
    \node at ($(opampA) + (0, -2)$) {\textbf{ADA4807-2} (A)};  

    %%%%%%%%%%%%%%%%%%%%%%%%%%%
    % Lower OpAmp (Channel B)
    %%%%%%%%%%%%%%%%%%%%%%%%%%%
    \draw (0,7) node[op amp] (opampB) {}
        % 1. Supply and ground
        (opampB.+) -- ++(-0.5,0) node[ground] {} 
        (opampB.up) -- ++(0,0.3) node[above] {+2.5V}
        (opampB.down) -- ++(0,-0.3) node[below] {-2.5V}
        
        % 2. Inverting input -> Point iouta -> Current source
        (opampB.-) to[short, -o] ++(-2,0) coordinate (ioutb) node[left] {\(I_{outB}\)}
        (ioutb) -- ++(0,-0.5) to[isource, l=$I_B$] ++(0,-1.5) node[ground]{}
        
        % 3. Back to inverting input and feedback
        (opampB.-) -- ++(0,2) coordinate (fbB) 
        to[R, l={\(R_{fb}\)}, a=249R] (fbB -| opampB.out) -- (opampB.out)
        
        % 4. Output
        (opampB.out) to[R, l={\(R_{b}\)}, a=88R7, -o] ++(2.5,0) node[right] {\(OUT_B\)}; 

    \node at ($(opampB) + (0, -2)$) {\textbf{ADA4807-2} (B)};  


    %%%%%%%%%%%%%%%%%%%%%%
    % ADA4870-1 OpAmp
    %%%%%%%%%%%%%%%%%%%%%%
    \draw (10,10) node[op amp] (pwrAmp) {};
    \node at ($(pwrAmp) + (0, -2)$) {\textbf{ADA4870-1}};

    % Non-inverting input (+) -> OUT_B
    \draw (pwrAmp.+) to[short, -o] ++(-1.0,0) node[left] {\(OUT_B\)};

    % Inverting input (-) -> OUT_A
    \draw (pwrAmp.-) to[short, -o] ++(-1.0,0) node[left] {\(OUT_A\)};

    % Feedback path (Rf || Cf)
    \draw (pwrAmp.-) -- ++(0,2) coordinate (fb_node) 
            node[circ] {}; % connection point at input
            
    \draw (fb_node) to[R, l={\(R_f\)}, a={1K4}] (fb_node -| pwrAmp.out) 
            coordinate (fb_end) node[circ] {} % connection point at output (Rf and output line)
            -- (pwrAmp.out) node[circ] {}; % point at opamp output

    % Capacitor Cf (parallel to Rf)
    \draw (fb_node) -- ++(0,2) coordinate (cap_node) 
            to[C, l={\(C_f\)}, a={1pF}] (cap_node -| fb_end) 
            -- (fb_end);

    % Supply
    \draw (pwrAmp.up) to[short, -o] ++(0,0.5) node[above] {\(V_{DD}\)};
    \draw (pwrAmp.down)  to[short, -o] ++(0,-0.5) node[below] {\(V_{EE}\)};

    % Output section
    \draw (pwrAmp.out) -- ++(1.5,0) coordinate (output_bus)
            node[circ] {}; % connection point to load

    % Load capacitor to ground
    \draw (output_bus) to[C, l=$C_{load}$, a={\(220pF\)}] ++(0,-2) node[ground] {};

    % Output label Signal
    \draw (output_bus) to[short, -o] ++(1.5,0) node[font=\boldmath, shift={(5pt,10pt)}] {Signal};

\end{circuitikz}
\caption{Overall connection diagram of ADA4807-2 to ADA4870-1}
\end{figure}

\begin{itemize}
    \item The differential pair \(I_{outA}\), \(I_{outB}\) of the AD9705 are represented by two current sources.
    \item The current sources generate about 4~mA (the maximum output current of the AD9705 per datasheet does not exceed 5~mA)
    \item Parasitic parameters of the current sources:
    \begin{itemize}
        \item Output capacitance of channels \(I_{outA}\), \(I_{outB}\) --- 5~pF
        \item Output resistance of channels \(I_{outA}\), \(I_{outB}\) --- 150~k\(\Omega\)
    \end{itemize}
    \item The signal of the first source is defined by:
    \begin{align}
        I=\frac{I_{FS}}{2}  + \frac{I_{FS}}{2} \left( \sin{2\pi \cdot \omega \cdot t} \right)
    \end{align}
    \item The signal of the second source is defined by (shifted by \(\pi\)):
    \begin{align}
        I=\frac{I_{FS}}{2}  + \frac{I_{FS}}{2} \left( \sin{2\pi \cdot \omega \cdot t  + \pi} \right)
    \end{align}
    \item \(\omega\) is the frequency. In this case, 1~MHz
    \item \(V_{DD} = 18V\),
    \item \(V_{EE} = -18V\)
\end{itemize}

\par In LTspice the schematic looks like:

\begin{figure}[H]
    \centering
    \includegraphics[width=0.9\linewidth, page=1]{scheme.png}
    \caption{LTspice simulation schematic}
\end{figure}

% =============================================
% 2. CIRCUIT DESCRIPTION
% =============================================
\section{Circuit Description}

\subsection{Block Diagram}

The circuit consists of the following functional blocks:
\begin{enumerate}
  \item Current sources \(I_{OUTA}\) and \(I_{OUTB}\) --- models of the DAC AD9705 current outputs
  \item Input transimpedance amplifiers (TIA) based on the dual ADA4807-2 (U2, U3), converting current to voltage
  \item ADA4870 input network: resistors \(R_{a}\), \(R_{b}\) and parasitic elements
        \(R_{para}\)/\(C_{para}\) (models of DAC output impedance)
  \item ADA4870 differential amplifier (U1) with feedback network \(R_f\)/\(C_f\)
  \item Load network: \(R_{Load}\) and \(C_{Load}\)
  \item Voltage divider \(R_1\)/\(R_2\) for SD input control
\end{enumerate}

\subsection{Current Sources \texorpdfstring{(\(I_{OUTA}\), \(I_{OUTB}\))}{{(IOUTA, IOUTB)}}}

The sources model the differential current output of the AD9705 DAC
\cite{AD9705_ds}. The currents vary in antiphase:
\begin{align}
  I_{\text{OUTA}} &= \frac{I_{FS}}{2}
    + \frac{I_{FS}}{2}\sin(2\pi f t),
    \label{eq:iouta}\\
  I_{\text{OUTB}} &= \frac{I_{FS}}{2}
    + \frac{I_{FS}}{2}\sin(2\pi f t + \pi),
    \label{eq:ioutb}
\end{align}
where $I_{FS}$ is the DAC full-scale current (in the current simulation $I_{FS}=4$~mA,
which corresponds to reduced output current operation),
$f = 1$~MHz is the signal frequency.

Each source is unidirectional: the current varies from 0 to $I_{FS}$, the DC component
$I_{FS}/2$ is cancelled by the differential configuration.
In the simulation, \texttt{BI} sources (LTspice behavioral current source) are used.

\subsection{Transimpedance Amplifiers (ADA4807-2)}

Two channels of the dual ADA4807-2 (U2, U3) operate as current-to-voltage
converters (TIA).
Each channel connection:
\begin{itemize}
    \item Non-inverting input connected to ground (GND).
    \item Inverting input connected to the DAC current output
          (\(I_{OUTA}\) for U2, \(I_{OUTB}\) for U3).
    \item Feedback resistor \(R_3\) (249~Ω) connected between output and
          inverting input (for U2); \(R_5\) (249~Ω) for U3.
\end{itemize}
The virtual ground at the inverting input ensures that the input current
flows through the feedback resistor, generating the output voltage:
\begin{align}
    V_{N001} &= -I_{OUTA} \cdot R_3, \label{eq:v_n001}\\
    V_{N002} &= -I_{OUTB} \cdot R_5. \label{eq:v_n002}
\end{align}

The differential voltage at the ADA4870 inputs (ASA and ASB) is
\begin{equation}
    V_{\text{diff}} = V_{N001} - V_{N002}
    = -(I_{OUTA} - I_{OUTB}) \cdot R_3
    = -I_{FS}\sin(2\pi f t) \cdot 249\,\Omega,
\end{equation}
since \(I_{OUTA} - I_{OUTB} = I_{FS}\sin(2\pi f t)\).
The amplitude of the differential signal is
\[
    V_{\text{diff,amp}} = I_{FS} \cdot R_3 = 0.004\,\text{A} \times 249\,\Omega
    = 0.996\,\text{V}.
\]

The ADA4807-2 is powered from sources V3, V4: $\pm 2.5$~V.
Resistors \(R_3\) and \(R_5\) are chosen as 249~Ω from the standard
series, ensuring high linearity of conversion at currents up to 4~mA.

\subsection{ADA4870 Input Network}

\subsubsection{Resistors \texorpdfstring{\(R_{a}\) and \(R_{b}\)}{Ra and Rb}}

\textbf{\(R_{a}\)} (150~Ω in the current revision) is connected
\emph{in series} between the output of U2 (node N001) and the inverting
input of ADA4870 (node OUTA). This node also carries the feedback
resistor \(R_f\) — so \(R_a\) is in the low-impedance summing node
circuit, through which real signal current flows (not a shunt to ground,
as one might assume by analogy with a classic 4-resistor differential
amplifier).

\textbf{\(R_{b}\)} (140~Ω) is connected \emph{in series} between the
output of U3 (node N002) and the non-inverting input of ADA4870 (node OUTB).
There is no free shunt to ground in this topology on either input.

The formula
\begin{equation}
  \boxed{R_b = \frac{R_a \cdot R_f}{R_a + R_f}}
  \label{eq:sym}
\end{equation}
remains valid, but its physical meaning is input bias current matching:
it equalizes the DC resistance seen by each op-amp input, thereby removing
the contribution of input bias current to output offset. \textbf{It does
not nullify the DC component at the output} — contrary to what the name
"symmetry condition" might suggest.

The DC component at the output arises from the very nature of the DAC currents:
\(I_{OUTA}\) and \(I_{OUTB}\) are unidirectional, their sum
\(I_{OUTA}+I_{OUTB}=I_{FS}=\text{const}\), not zero. This common-mode
current flows through the asymmetric (single feedback resistor on the
inverting side) ADA4870 network and produces an output offset
\begin{equation}
  V_{\text{offset}} = -\frac{R_{TIA}\cdot I_{FS}}{2},
  \label{eq:offset}
\end{equation}
which \textbf{does not depend on the choice of \(R_a\), \(R_f\), \(R_b\)} —
only on \(R_{TIA}\) and \(I_{FS}\). With \(R_{TIA}=249\)~Ω,
\(I_{FS}=4\)~mA: \(V_{\text{offset}} \approx -0.498\)~V — this is exactly
the \(\sim{-0.5}\)~V observed in simulation, not a consequence of
inaccurate E96 rounding of \(R_b\) (see also section~\ref{sec:algorithm}
and the corrected note in section~\ref{sec:fourier820}).

Separately from the offset — the gain asymmetry between the rails
(see equation~\eqref{eq:av} and the comment after it): this is also
structural, but a different effect not directly related to offset,
and also not fixable by tuning \(R_a\)/\(R_b\).

\textbf{Important:} swapping \(R_{a}\) and \(R_{b}\) would violate
condition~\eqref{eq:sym}, because node OUTA additionally has the feedback
resistor \(R_f\). Therefore \(R_b\) must always equal \(R_a \parallel R_f\),
not the reverse.

\subsubsection{DAC Parasitic Elements (\texorpdfstring{\(R_{paraA}\)/\(C_{paraA}\), \(R_{paraB}\)/\(C_{paraB}\)}{Rpara/Cpara})}

These elements model the output impedance of the AD9705 current outputs.
Typical values from the datasheet \cite{AD9705_ds} are used:
\begin{itemize}
    \item \(R_{paraA} = R_{paraB} = 150\)~kΩ,
    \item \(C_{paraA} = C_{paraB} = 5\)~pF.
\end{itemize}
In parallel, capacitors Ca = Cb = 22~pF are also included (model of parasitic
board capacitance). Thus, total capacitance at each input node (INN and INP)
is \(5+22 = 27\)~pF.

The influence at low frequencies (1~MHz) is negligible. The pole formed
by 27~pF and resistance of about 100~Ω lies at about 59~MHz, above the
operating band.

\subsection{Feedback Resistors \texorpdfstring{\(R_{fa}\), \(R_{fb}\)}{Rfa, Rfb}}
\par The value of 249 Ω for resistors \(R_{fa}\) and \(R_{fb}\) in the transimpedance amplifiers (TIA) on ADA4807-2 was not chosen arbitrarily, but based on a balance between several key requirements. Here is the rationale:

\begin{itemize}
    \item The DAC full-scale current was set to \(I_{FS} = 4 mA\). This current determines the signal amplitude at the TIA output.
    \item Target output voltage: To effectively drive the subsequent ADA4870 stage, which requires about 1 V differential to obtain ±14 V amplitude, a value \(V_{diff_amp} \approx 0.996\) V was chosen.
    \item For a classic TIA with op-amp, where current flows into the virtual ground of the inverting input, the output voltage is simply: \(V_{out} = I_{in} \times R_f\). In our case:
\begin{align}
    R_f = \frac{V_{diff\_amp}}{I_{FS}} = \frac{0.996V}{0.004A} \approx 249 \Omega
\end{align}    
\end{itemize}

\par Thus, 249 Ω is a standard E96 series value that perfectly converts 4 mA current into about 1 V.
\par A different value could have been used, but 249 Ω proved optimal because:
\par It does not overload the DAC: The DAC output voltage (about 1 V) remains within its compliance range, ensuring linearity.
\begin{itemize}
    \item It does not overload the ADA4807 input: The TIA output stays within the linear range of the op-amp with \(\pm 2.5V\) supply.
    \item It leaves gain margin: Even a small change in this resistor would require recalculation of \(Ra\) and \(R_b\) in the output stage to maintain balance.
\end{itemize}

\par If the DAC output current \(I_{FS}\) needs to be changed, this resistor should be recalculated using the same formula to keep the same differential voltage for the ADA4870 stage.

\subsection{Amplifier U1 (ADA4870)}

The ADA4870 is connected in a differential configuration.
The signal from U2 output (node N001) is fed to the \textbf{inverting}
input via resistor \(R_a\) (the same node OUTA carries the feedback
resistor \(R_f\)), while the signal from U3 output (node N002) is fed to
the \textbf{non-inverting} input via \(R_b\). Feedback is provided by
resistor \(R_f\) from the SIGNAL output to the inverting input (node OUTA).

The differential signal gain is:
\begin{equation}
    A_v = \frac{R_f}{R_a}.
    \label{eq:av}
\end{equation}

With \(R_f = 1.4\)~kΩ and \(R_a = 100\)~Ω, \(A_v = 14\).

\textbf{Note on asymmetry.} Formula~\eqref{eq:av} is the gain only for
the inverting path. For the non-inverting path (signal via \(R_b\)), the gain
is \(1+R_f/R_a\), which is exactly 1 larger in magnitude — this is a
topological property of the single-feedback-resistor circuit, independent
of the specific values. With target \(A_v\approx14\), the discrepancy
(14 vs 15) gives an output peak asymmetry of about 6--7\%, which is
considered acceptable given the available voltage headroom (see also
section~\ref{sec:algorithm}).

The output voltage amplitude for the given DAC current:
\begin{equation}
    V_{\text{peak}} = A_v \cdot V_{\text{diff,amp}}
    = \frac{R_f}{R_a} \cdot I_{FS} \cdot R_3.
    \label{eq:vpeak}
\end{equation}

For \(I_{FS}=4\)~mA, \(R_3=249\)~Ω, \(R_a=100\)~Ω, \(R_f=1.4\)~kΩ,
we obtain:
\[
    V_{\text{peak}} = 14 \times 0.996 \approx 13.94\ \text{V}.
\]
Peak load current for \(R_{Load}=22\)~Ω:
\[
    I_{\text{peak}} = \frac{14\ \text{V}}{22\ \Omega} \approx 0.636\ \text{A} < 1\ \text{A}.
\]

ADA4870 supply: $\pm 18$~V from sources V1, V2.

\subsection{Feedback Network \texorpdfstring{(\(R_f\), \(C_f\))}{Rf, Cf}}

\textbf{\(R_f\)} (1.4~kΩ) is the main feedback resistor defining the gain.

\textbf{\(C_f\)} (2~pF) is connected in parallel with \(R_f\). It introduces
a zero in the loop transfer function, compensating the pole from the
parasitic capacitance at the ASA input.
The zero frequency:
\[
    f_{z} = \frac{1}{2\pi R_f C_f}
    = \frac{1}{2\pi \times 1400 \times 2\times 10^{-12}}
    \approx 56.8\ \text{MHz}.
\]
The optimal value is \(C_f = C_{\text{INN}} \cdot R_a / R_f
\approx 27\ \text{pF} \times 100 / 1400 \approx 1.9\ \text{pF}\).
The used value of 2~pF is close to the calculated one and ensures stability.

\subsection{Load Network \texorpdfstring{(\(R_{Load}\), \(C_{Load}\))}{Rload, Cload}}

\textbf{\(R_{Load}\)} = 22~Ω — active load resistance.
Peak current 0.636~A, well below the maximum ADA4870 current (1~A).

\textbf{\(C_{Load}\)} = 220~pF — model of parasitic board and load capacitance.
At 1~MHz its impedance is about 723~Ω, much larger than 22~Ω, so the
shunting effect is minimal.

\subsection{SD Divider (R1, R2)}

The divider formed by R1 (2~kΩ) and R2 (10~kΩ) sets the voltage on the
SD enable input of the ADA4870.
With $\pm 18$~V supply, the SD voltage is
\(-18\ \text{V} \times \frac{10}{12} = -15\ \text{V}\),
which falls within the enable range (from \(-16.9\ \text{V}\)
to \(-13\ \text{V}\)).

% =============================================
% 3. THEORETICAL CALCULATION
% =============================================
\section{Theoretical Calculation}

\subsection{Component Selection Algorithm}
\label{sec:algorithm}

Based on the given requirements (\(V_{out\_amp}\), \(R_{load}\), \(I_{FS}\)) and
component datasheets, an automated algorithm for selecting \(R_a, R_f, R_b, C_f\)
has been developed. Key steps:

\begin{enumerate}
  \item \textbf{Feasibility check.}
    \begin{itemize}
      \item Maximum undistorted amplitude: \(V_{max} = |V_{sup}| - V_{headroom}\).
      \item Peak load current: \(I_{peak} = V_{out\_amp}/R_{load}\). If \(I_{peak} > I_{out\_max}\) or \(V_{out\_amp} > V_{max}\), calculation stops.
    \end{itemize}
  \item \textbf{TIA differential voltage.}
    \(V_{diff\_amp} = I_{FS} \cdot R_{TIA}\) (typically 0.996 V for 4 mA and 249 Ω).
  \item \textbf{Required gain.}
    \(A_v = V_{out\_amp} / V_{diff\_amp}\).
  \item \textbf{Iterate over \(R_a\) candidates.}
    For each value from the list \texttt{Ra\_candidates}:
    \begin{align*}
      R_f &= A_v \cdot R_a;\\
      R_b &= \frac{R_a \cdot R_f}{R_a + R_f}\quad \text{(ideal balance condition)};\\
      C_f &= C_{f\_base} + \frac{\max(0, R_f - R_{f\_threshold})}{1000} \cdot C_{f\_per\_kohm}.
    \end{align*}
    \(R_f\) and \(R_b\) are rounded to the nearest E96 (1\%) series value.
    Variants with \(R_f > R_{f\_max}\) are discarded.
  \item \textbf{Ranking variants.}
    All valid combinations are sorted by the criterion of \textbf{minimum balancing error of \(R_b\)}:
    \[
      \delta R_b = |R_{b\_ideal} - R_{b\_E96}| \to \min.
    \]
    The smaller \(\delta R_b\), the lower the DC component at the ADA4870 output. In case of equal \(\delta R_b\), preference is given to the combination where \(R_f\) is closest to the manufacturer recommended value (see Table~\ref{tab:rf_rec}).
\end{enumerate}

\subsection{Recommended \(R_f\) Values for ADA4870}
\label{sec:rf_recommendation}

The ADA4870 is a current-feedback amplifier (CFB), and the choice of \(R_f\)
significantly affects bandwidth and stability. The manufacturer recommends
optimal \(R_f\) for typical gains (see Table 6 in the datasheet~\cite{ADA4870_ds}):

\begin{table}[H]
\centering
\caption{Recommended \(R_f\) for ADA4870}
\label{tab:rf_rec}
\begin{tabular}{@{}cc@{}}
\toprule
\textbf{Gain (V/V)} & \textbf{\(R_f\) (Ω)} \\
\midrule
+1  & 2000 \\
-1  & 1210 \\
+2  & 1500 \\
-2  & 1210 \\
+5, +10 & 1210 \\
\bottomrule
\end{tabular}
\end{table}

In our differential configuration, the effective gain exceeds 10, so the target
guideline for \(R_f\) is \textbf{1.21 kΩ}. This guideline is used as a secondary
sorting key.

\subsection{Example Calculation for this Circuit}
\label{sec:example_calc}

\begin{itemize}
    \item \(V_{\text{peak}} = 14\)~V, \(R_{Load} = 22\)~Ω $\Rightarrow I_{\text{peak}} = 0.636$~A.
    \item \(I_{FS} = 4\)~mA, \(R_{TIA} = 249\)~Ω $\Rightarrow V_{\text{diff,amp}} = 0.996$~V.
    \item \(R_a = 100\)~Ω, required \(A_v = 14 / 0.996 \approx 14.06\).
    \item \(R_f = 14.06 \times 100 = 1406\)~Ω $\to$ choose 1.4~kΩ (E96: 1.40) or 1.43~kΩ (closer).
        We take \(R_f = 1400\)~Ω, then \(A_v = 14.00\).
    \item \(R_b = (100 \times 1400)/(1500) = 93.33\)~Ω $\to$ 93.1~Ω (E96).
    \item \(C_f \approx (27\times10^{-12}) \times 100 / 1400 \approx 1.9\)~pF $\to$ keep 2~pF.
    \item Check: \(V_{\text{peak}} = A_v \cdot V_{\text{diff,amp}} = 14 \times 0.996 = 13.94\)~V,
        meeting the requirement.
\end{itemize}

\subsection{Final Component Values}

\textbf{Updated:} the table below reflects the current algorithm selection
(section~\ref{sec:algorithm}) with current \texttt{Ra\_candidates} and
\texttt{Rf\_target=1210} in the config — \(R_a=150\)~Ω,
\(R_f=2100\)~Ω (E96), \(R_b=140\)~Ω (E96, exact match to
\(R_a\parallel R_f\), \(\delta R_b=0\)). The example in section~\ref{sec:example_calc}
above (with \(R_a=100\)~Ω) is an earlier iteration, kept for illustration of
the calculation flow, but is not the current operating point.

\begin{table}[H]
\centering
\caption{Component values of schematic ada4807\_4870.asc (current revision)}
\label{tab:components}
\begin{tabular}{@{}llp{6.5cm}@{}}
\toprule
\textbf{Component} & \textbf{Value} & \textbf{Purpose} \\
\midrule
\(R_f\)     & 2.1~kΩ & Transimpedance, sets amplitude \\
\(R_a\)     & 150~Ω   & Resistor to inverting input of U1 \\
\(R_b\)     & 140~Ω & Resistor to non-inverting input of U1 (\(R_a \parallel R_f\)) \\
\(R_3, R_5\) & 249~Ω   & TIA resistors (current-to-voltage conversion) \\
\(C_f\)     & 1.9~pF     & Feedback compensation \\
Ca, Cb     & 22~pF    & Parasitic capacitance of DAC inputs \\
\(R_{paraA}, R_{paraB}\) & 150~kΩ & DAC output resistance \\
\(C_{paraA}, C_{paraB}\) & 5~pF   & DAC output capacitance \\
\(C_{Load}\)  & 220~pF   & Parasitic load capacitance \\
\(R_{Load}\)  & 22~Ω    & Load \\
R1, R2     & 2~kΩ, 10~kΩ & SD voltage divider \\
\(V_{sup}\) & $\pm 18$~V & ADA4870 supply \\
\(V_{sup,TIA}\) & $\pm 2.5$~V & ADA4807-2 supply \\
\bottomrule
\end{tabular}
\end{table}

\subsection{Adaptive Simulation Time Setting}
\label{sec:adaptive_tran}

To obtain reliable Fourier analysis results at any frequency, it is necessary
that an integer number of signal periods fits into the analysis window and
that transients have settled. The program implements dynamic generation of
the \texttt{.tran} directive based on parameters set in \texttt{config.yaml}
(previously \texttt{config.json} before switching to YAML):

\begin{itemize}
    \item \texttt{periods\_transient} = 10 – number of periods for transient decay;
    \item \texttt{periods\_analysis} = 10 – number of periods whose data are saved for analysis;
    \item \texttt{points\_per\_period} = 1000 – minimum number of points per period.
\end{itemize}

From these values, the following are automatically computed:
\begin{align}
    T &= \frac{1}{f_{\text{signal}}},\\
    t_{\text{end}} &= (N_{\text{transient}} + N_{\text{analysis}}) \cdot T,\\
    t_{\text{start}} &= N_{\text{transient}} \cdot T,\\
    \Delta t_{\text{max}} &= \frac{T}{N_{\text{points}}}.
\end{align}

Thus, for low frequencies the simulation becomes longer, but guarantees that
transients have settled before saving starts, and Fourier analysis is performed
on an integer number of periods. This eliminates spurious THD behavior and
ensures a physically meaningful frequency dependence of distortion.

% =============================================
% 4. THERMAL CONDITIONS
% =============================================
\section{Thermal Conditions of ADA4870}

The average power dissipated by the amplifier with a sinusoidal
signal \cite{ADA4870_ds}:
\begin{equation}
  P_{\text{AVG}} = V_S \cdot I_q
    + \frac{2}{\pi}\cdot\frac{V_{CC}\,V_{\text{peak}}}{R_L}
    - \frac{V_{\text{peak}}^2}{2R_L},
  \label{eq:power}
\end{equation}
where \(V_S = 36\)~V (total), \(I_q = 32.5\)~mA,
\(V_{CC} = 18\)~V, \(V_{\text{peak}} = 14\)~V, \(R_L = 22\)~Ω:
\[
    P_{\text{AVG}} = (36 \times 0.0325)
    + \frac{2}{\pi}\cdot\frac{18 \times 14}{22}
    - \frac{14^2}{2 \times 22}
    = 1.17 + \frac{2}{\pi}\cdot 11.4545 - 4.4545
    \approx 1.17 + 7.29 - 4.45 = 4.01\ \text{W}.
\]

Junction temperature at ambient \(T_{\text{amb}} = 40\)~°C
and case-to-ambient thermal resistance \(\theta_{CA}\):
\[
    T_J = T_{\text{amb}} + P_{\text{AVG}} \cdot (\theta_{JC} + \theta_{CA}),
\]
where \(\theta_{JC} = 1.3\)~°C/W (for SOIC package).
For a safe margin (\(T_J \le 150\)~°C), we require
\[
    \theta_{CA} \le \frac{150 - 40}{4.01} - 1.3 \approx 27.4 - 1.3 = 26.1\ \text{°C/W}.
\]
A heatsink with thermal resistance not exceeding 20~°C/W is recommended.

Load power:
\[
    P_{R_{\text{load}}} = \frac{V_{\text{peak}}^2}{2R_{\text{load}}}
    = \frac{196}{44} \approx 4.45\ \text{W}.
\]
The load resistor should be rated for at least 10~W.

% =============================================
% 5. SIMULATION RESULTS
% =============================================
\section{LTspice Simulation Results}

Simulation was performed in LTspice version 26.0.1 for Windows \cite{LTspice}.
\textbf{Note:} the run below is an early iteration with \(R_f = 820\)~Ω
(historical, for illustration of the methodology), giving an amplitude of
about \(8.65\)~V on a 22~Ω load, which is lower than the target 14~V.
Current values and results are in section~\ref{sec:degradation} and
Table~\ref{tab:components}.
Simulation directives:
\begin{itemize}
  \item \texttt{.tran 10us}
  \item \texttt{.FOUR 1Meg V(aout)}
  \item \texttt{.MEAS TRAN Vfinal FIND V(aout) AT=10us}
\end{itemize}

\subsection{Simulation Parameters}

\begin{table}[H]
\centering
\caption{Simulation parameters (.param)}
\label{tab:params}
\begin{tabular}{@{}lll@{}}
\toprule
\textbf{Parameter} & \textbf{Value} & \textbf{Purpose} \\
\midrule
IFS   & 4~mA   & DAC full-scale current \\
FREQ  & 1~MHz  & Signal frequency \\
\(R_f\)Val & 820~Ω & Feedback resistor (base configuration) \\
\(R_{Load}\)Val & 22~Ω & Load resistance \\
\bottomrule
\end{tabular}
\end{table}

\subsection{Fourier Analysis of Output Signal (Rf = 820 Ω)}
\label{sec:fourier820}

\begin{table}[H]
\centering
\caption{Fourier analysis of V(aout), \(f_0 = 1\)~MHz (from .log file)}
\label{tab:fourier}
\begin{tabular}{@{}rrrrrr@{}}
\toprule
\textbf{No.} &
\textbf{Frequency, Hz} &
\textbf{Amplitude, V} &
\textbf{Norm.} &
\textbf{Phase, °} &
\textbf{Norm. phase, °} \\
\midrule
1 & \(1.000\times10^6\) & \(8.645\times10^0\) & \(1.000\) & \(92.73\) & \(0.00\) \\
2 & \(2.000\times10^6\) & \(1.647\times10^{-4}\) & \(1.905\times10^{-5}\) & \(140.79\) & \(48.06\) \\
3 & \(3.000\times10^6\) & \(1.254\times10^{-4}\) & \(1.450\times10^{-5}\) & \(-69.62\) & \(-162.35\) \\
4 & \(4.000\times10^6\) & \(1.373\times10^{-4}\) & \(1.588\times10^{-5}\) & \(94.03\) & \(1.30\) \\
5 & \(5.000\times10^6\) & \(1.213\times10^{-4}\) & \(1.403\times10^{-5}\) & \(-127.85\) & \(-220.57\) \\
6 & \(6.000\times10^6\) & \(1.266\times10^{-4}\) & \(1.465\times10^{-5}\) & \(37.79\) & \(-54.94\) \\
7 & \(7.000\times10^6\) & \(1.439\times10^{-4}\) & \(1.665\times10^{-5}\) & \(179.05\) & \(86.32\) \\
8 & \(8.000\times10^6\) & \(1.378\times10^{-4}\) & \(1.594\times10^{-5}\) & \(-13.93\) & \(-106.65\) \\
9 & \(9.000\times10^6\) & \(1.650\times10^{-4}\) & \(1.908\times10^{-5}\) & \(137.08\) & \(44.35\) \\
\bottomrule
\end{tabular}
\end{table}

\begin{table}[H]
\centering
\caption{Simulation summary (Rf = 820 Ω)}
\label{tab:summary}
\begin{tabular}{@{}lr@{}}
\toprule
\textbf{Parameter} & \textbf{Value} \\
\midrule
Fundamental amplitude       & 8.645~V \\
DC offset                  & \(-0.5446\)~V \\
Total THD                    & 0.00633\% \\
\(V_{\text{aout}}\) at \(t=10\)~µs & \(-0.956\)~V \\
\(V_{SD}\)                      & \(-15.183\)~V \\
\bottomrule
\end{tabular}
\end{table}

\textbf{Corrected:} the measured DC shift (\(-0.54\) V) is not a consequence
of rounding \(R_b\) to 88.7~Ω instead of the exact 93.33~Ω (this error
contributes two orders of magnitude less). It is the structural offset
\(V_{\text{offset}} = -R_{TIA}\cdot I_{FS}/2 \approx -0.498\)~V
(see equation~\eqref{eq:offset}) — it does not depend on \(R_b\) and is not
eliminated by precise tuning of \(R_b\), whether with \(R_f=820\)~Ω or
\(R_f=1.4\)~kΩ.

\subsection{Expected Performance with Rf = 1.4 kΩ}

With \(R_f\) increased to 1.4~kΩ and the circuit balanced, the following
parameters are expected:
\[
    V_{\text{peak}} \approx 14\ \text{V},\quad
    \text{DC offset} \approx -0.5\ \text{V},\quad
    \text{THD} < 0.01\%.
\]
Peak current into 22~Ω load will be about 0.64~A.
The ADA4870 power dissipation will increase to \(\sim 4\)~W,
requiring a heatsink.

\section{Simulation Results}

\par At the outputs \(I_{NN}\) and \(I_{NP}\) we obtain:
\begin{figure}[H]
\centering
\begin{tikzpicture}
\begin{axis}[
    width=12cm, height=7cm,
    xlabel={Time (s)},
    ylabel={Voltage \(V(inn)\) (V)},
    grid=major,
    legend pos=north east,
    table/col sep=comma 
]
    % Plot for V(inn)
    \addplot[red, thick] table [
        x=Time,
        y=V(inn),
    ] {./data/ada4870_raw_export.csv};
    \addlegendentry{$V(inn)$}

    % Plot for V(inp)
    \addplot[blue, thick] table [
        x=Time,
        y=V(inp)
    ] {./data/ada4870_raw_export.csv};
    \addlegendentry{$V(inp)$}
    
    \legend{\(V_{inn}\), \(V_{inp}\)}
\end{axis}
\end{tikzpicture}
\caption{Voltage at nodes \(I_{NN}\), \(I_{NP}\)}
\end{figure}

\par Then, the output signal looks like:

\begin{figure}[H]
\centering
\begin{tikzpicture}
\begin{axis}[
    width=12cm, height=7cm,
    xlabel={Time (s)},
    ylabel={Voltage \(V_{signal}\) (V)},
    grid=major,
    legend pos=north east,
    table/col sep=comma 
]
    % Plot for V(signal)
    \addplot[green!50!black, thick] table [
        x=Time,
        y=V(signal)] {./data/ada4870_raw_export.csv};

    \addlegendentry{$V_{signal}$}
   
    \legend{\(V_{signal}\)}
\end{axis}
\end{tikzpicture}
\caption{ADA4870-1 output voltage \(V_{signal}\)}
\end{figure}

\begin{figure}[H]
\centering
\begin{tikzpicture}
\begin{axis}[
    width=12cm, height=7cm,
    xlabel={Time (s)},
    ylabel={Current through \(R_{load}\) (A)},
    grid=major,
    legend pos=north east,
    table/col sep=comma 
]
    % Plot for I(load)
    \addplot[brown, thick] table [
        x=Time,
        y=I(Rload)] {./data/ada4870_raw_export.csv};

    \legend{\(I_{load}\)}
\end{axis}
\end{tikzpicture}
\caption{Current through \(R_{load}\)}
\end{figure}

\subsection{THD Degradation vs Frequency}
\label{sec:degradation}

To evaluate the operating frequency range and the behaviour of nonlinear
distortion, the \texttt{degradation\_sweep} method is used – sequential
simulation runs at several frequencies with fixed component values. To avoid
artefacts related to transients and insufficient analysis duration, adaptive
simulation time settings are used (see section~\ref{sec:adaptive_tran}).

Figure~\ref{fig:thd_vs_freq} shows the THD vs frequency plot obtained in the
range 100\,kHz – 5\,MHz for the selected component combination (\(R_a=150\)\,Ω,
\(R_f=2100\)\,Ω, \(R_b=140\)\,Ω). \textbf{Corrected:} the dependence is not
monotonic — THD rises from \(\sim 0.060\%\) (100\,kHz) to a peak of
\(\sim 0.063\%\) around 500\,kHz–1\,MHz, then decreases to \(\sim 0.049\%\) at
5\,MHz. This behaviour is atypical for simple reduction of feedback depth with
frequency (which would give monotonic increase) — a possible reason is that
\(C_f\) together with \(R_f\) creates a pole that suppresses higher harmonics
more strongly at the upper end of the range than the amplifier's own distortion
increases; this requires separate verification if the exact physical mechanism
is needed.

\begin{figure}[H]
\centering
\begin{tikzpicture}
\begin{axis}[
    width=12cm, height=7cm,
    xlabel={Frequency (Hz)},
    ylabel={THD (\%)},
    grid=major,
    legend pos=north west,
    xmode=log,
    log basis x={10},
    table/col sep=comma
]
    % Data from CSV generated by degradation_sweep
    \addplot[blue, thick, mark=*] table [
        x=Frequency_Hz,
        y index=1
    ] {./data/thd_vs_freq.csv};
    \legend{THD}
\end{axis}
\end{tikzpicture}
\caption{THD degradation vs frequency. Data obtained via automated sweeping.}
\label{fig:thd_vs_freq}
\end{figure}

The obtained dependence confirms that throughout the range 100\,kHz–5\,MHz,
distortion remains very low (THD $\lesssim 0.063\%$, with a small peak in the
middle rather than monotonically increasing towards the upper bound) at an
output power of $\sim 4.5$~W. The distortion margin for the target frequencies
(up to 2~MHz per specification) is more than sufficient.

% =============================================
% 6. RECOMMENDATIONS
% =============================================
\section{Output Protection Recommendations}

\subsection{ADA4870 Output Protection}

When driving inductive loads, voltage spikes beyond the supply rails may occur.
It is recommended to install Schottky protection diodes directly at the AOUT pin.

Connection diagram (Fig.~\ref{fig:diode_scheme}):
\begin{figure}[H]
\centering
\begin{circuitikz}[american, scale=1.1]
    \draw (0,0) node[op amp, noinv input up, label={[yshift=-1cm, xshift=-0.5cm]below:ADA4870}] (oa) {};
    \draw (oa.out) -- ++(0.5,0) coordinate (out_node);
    \draw (oa.+) node[ocirc] {};
    \draw (oa.-) node[ocirc] {};
    \draw (out_node) to[Do, l_=VD1, a^={MBRS360}, *-*] ++(0, 2)
          node[above] {$+18$~V (VDD)};
    \draw (out_node) to[Do, l_=VD2, a^={MBRS360}, *-*] ++(0,-2)
          node[below] {$-18$~V (VEE)};
    \draw (out_node) to[R, l=$R_{\text{Load}}$\ 22~Ω, *-o] ++(3,0)
          node[right] {$V_{\text{out}}$};
\end{circuitikz}
\caption{Protection diode connection at ADA4870 output.}
\label{fig:diode_scheme}
\end{figure}

\begin{itemize}
  \item \textbf{VD1} (to positive rail): anode to AOUT,
        cathode to +18~V (VDD).
  \item \textbf{VD2} (to negative rail): anode to $-18$~V
        (VEE), cathode to AOUT.
\end{itemize}

\subsubsection{Diode Selection MBRS360}
MBRS360 \cite{MBRS360_ds} is a 3~A, 60~V Schottky diode in SMC package.
The peak load current of 0.636~A is well below the limit, reverse voltage
60~V safely covers possible spikes. The diode capacitance (about 180~pF) has
negligible effect on signals up to 2~MHz.

\subsection{Components for the Actual PCB}

\begin{itemize}
  \item \textbf{\(R_{Load}\)}: non-inductive resistor 22~Ω,
        rated for at least 10~W.
  \item \textbf{Protection diodes}: MBRS360 (or similar with current rating \(\ge 1\)~A).
  \item \textbf{Bypass capacitors}: 100~nF + 10~µF
        from each supply rail to GND, placed close to ADA4870.
  \item \textbf{\(C_f\)}: start with 2~pF, if necessary
        increase to 2.7–3.3~pF based on AC analysis.
  \item \textbf{TIA circuit}: resistors \(R_3=R_5=249\)~Ω,
        no capacitors in the ADA4807-2 feedback loop are required
        (driving a low-impedance load).
\end{itemize}

% =============================================
% BIBLIOGRAPHY
% =============================================
\begin{thebibliography}{99}

\bibitem{ADA4870_ds}
Analog Devices, Inc. \textit{ADA4870: High Speed, High Voltage,
1\,A Output Drive Amplifier. Data Sheet.} Rev.\,C, 2025.
URL: \url{https://www.analog.com/en/products/ada4870.html}.

\bibitem{ADA4807_ds}
Analog Devices, Inc. \textit{ADA4807-1/ADA4807-2/ADA4807-4:
Low Noise, Rail-to-Rail Amplifier. Data Sheet.} Rev.\,D, 2020.
URL: \url{https://www.analog.com}.

\bibitem{AD9705_ds}
Analog Devices, Inc. \textit{AD9704/AD9705/AD9706/AD9707:
8-/10-/12-/14-Bit, 175\,MSPS TxDAC Digital-to-Analog Converters.
Data Sheet.} Rev.\,G, 2023.
URL: \url{https://www.analog.com}.

\bibitem{MBRS360_ds}
ON Semiconductor. \textit{MBRS360T3G: Surface Mount Schottky Power
Rectifier. Data Sheet.} Rev.\,14, March 2020.
URL: \url{https://www.onsemi.com}.

\bibitem{LTspice}
Analog Devices, Inc. \textit{LTspice Simulation Software.}
Version~26.0.1, 2026.

\end{thebibliography}

\end{document}
